\documentclass[12pt,a4paper]{article}

\usepackage[utf8]{inputenc}
\usepackage[T1]{fontenc}
\usepackage[spanish]{babel}
\usepackage{amsmath,amssymb}
\usepackage{geometry}
\usepackage{graphicx}
\usepackage{float}
\usepackage{booktabs}
\usepackage{array}

\geometry{margin=2.5cm}

\title{Esquema de Esteganografía y Marca de Agua Robusta\\ Basado en ANS y STC\\[4pt]
\large Fundamentos Matemáticos y Teóricos del Sistema Implementado}
\author{
Michael S. Betancourt Gelves \\
Nicolás A. Castellanos Rico \\[6pt]
\small Profesor: Oswaldo Rojas Camacho \\
\small Teoría de la Información --- Universidad Nacional de Colombia, 2026
}
\date{}

\newcommand{\vect}[1]{\mathbf{#1}}

\begin{document}

\maketitle

\begin{abstract}
\noindent
Este documento formaliza los fundamentos matemáticos del sistema construido para
la protección y persistencia de marcas de agua invisibles. El sistema integra
\emph{dos modos complementarios}: (i)~una esteganografía de \emph{alta capacidad}
que combina compresión entrópica por Sistemas Numéricos Asimétricos (ANS),
funciones de costo adaptativas (J-UNIWARD) y Códigos de Síndrome-Trellis (STC)
sobre la Transformada Discreta del Coseno; y (ii)~una \emph{marca de agua robusta}
resistente al reescalado y la recompresión de las redes sociales, basada en
Espectro Ensanchado Mejorado (ISS) sobre una malla canónica, protección
Reed-Solomon, compresión entrópica del texto y una amplitud adaptativa por carga
útil. Para cada concepto matemático empleado se enuncia su formulación y se cita
su fuente bibliográfica original. El resultado logrado es una marca que sobrevive
la cadena completa de Facebook, WhatsApp y Pinterest (reescalado + recompresión
JPEG/WebP) conservando la imperceptibilidad.
\end{abstract}

\section*{1. Introducción: el desafío de la veracidad}

La irrupción de la Inteligencia Artificial generativa ha inundado la red de
\emph{deepfakes}. En este contexto ya no basta con \emph{ver para creer}: se
requieren mecanismos matemáticos para verificar la \emph{autenticidad} y la
\emph{integridad} de una imagen. Una marca de agua invisible de procedencia
(\emph{provenance}) permite incrustar una firma verificable dentro del propio
medio; el reto es que esa firma \emph{persista} cuando la imagen atraviesa canales
con pérdida.

El problema central es, por tanto, de \emph{Teoría de la Información}: transmitir
un mensaje oculto por un canal ruidoso (la compresión y el reescalado de las
plataformas) de forma imperceptible y recuperable. Este trabajo aborda ese
problema con herramientas clásicas de codificación de fuente, codificación de
canal, procesamiento de señales y modelado perceptual.

\section*{2. El problema: fragilidad ante canales con pérdida}

La esteganografía tradicional oculta bits en el dominio espacial o en bloques
$8\times 8$ de la Transformada Discreta del Coseno (DCT). Al subir una imagen a
una red social, algoritmos como JPEG la \emph{recomprimen} y, sobre todo, la
\emph{reescalan}: ambas operaciones desincronizan la retícula de $8\times 8$ y
destruyen los bits ocultos, produciendo un fallo de sincronización en la
extracción. El mensaje desaparece ante la más mínima edición.

Formalmente, si $x$ es el medio portador y $C(\cdot)$ el operador de canal
(reescalado seguido de recompresión), se necesita un esquema de codificación
$E$ y decodificación $D$ tales que
\[
D\big(C(E(x,m))\big) = m
\]
con alta probabilidad, manteniendo una distorsión perceptual pequeña entre $x$ y
$E(x,m)$. Las dos secciones finales (Modos A y B) presentan dos soluciones a este
problema con distintos compromisos capacidad/robustez.

\section*{3. Fundamentos de Teoría de la Información}

\subsection*{3.1 Entropía de Shannon y codificación de fuente}

La entropía de una fuente discreta $X$ con distribución $p(x)$ es
\begin{equation}
H(X) = -\sum_{x} p(x)\,\log_2 p(x),
\end{equation}
y cuantifica el número medio de bits por símbolo necesarios para representar la
fuente sin pérdida. El \emph{teorema de codificación de fuente} de Shannon
\cite{shannon1948} establece que ninguna codificación libre de pérdida puede usar
menos de $H(X)$ bits/símbolo en promedio, y que este límite es alcanzable
asintóticamente. Comprimir el mensaje \emph{antes} de incrustarlo reduce el número
de bits a ocultar; como se verá, menos bits implican menos modificaciones del
portador (menor distorsión) y mayor margen frente al canal.

\subsection*{3.2 Codificación entrópica: Huffman, Lempel-Ziv y ANS}

Para acercarse al límite $H(X)$ el sistema emplea codificación entrópica. Tres
resultados matemáticos la sustentan:

\begin{itemize}
\item \textbf{Códigos de Huffman} \cite{huffman1952}: construyen un código prefijo
de \emph{redundancia mínima} para una distribución conocida, óptimo entre los
códigos de longitud entera.
\item \textbf{Codificación por diccionario Lempel-Ziv (LZ77)} \cite{zivlempel1977}:
sustituye subcadenas repetidas por referencias $(\text{distancia},\text{longitud})$,
capturando la redundancia \emph{contextual} del lenguaje natural sin conocer la
distribución de antemano.
\item \textbf{Sistemas Numéricos Asimétricos (ANS)} \cite{duda2014}: representan el
mensaje como un único número natural cuyo estado $x$ evoluciona, para un símbolo
$s$ de frecuencia $f_s$ y frecuencia acumulada $c_s$ sobre un total
$M=2^{b}$, según
\begin{equation}
x' \;=\; \Big\lfloor \tfrac{x}{f_s} \Big\rfloor\, M \;+\; \big(x \bmod f_s\big) \;+\; c_s .
\end{equation}
ANS combina la \emph{velocidad} de Huffman con la \emph{eficiencia} de la
codificación aritmética \cite{rissanen1979}, situándose muy cerca del límite de
entropía.
\end{itemize}

En el Modo~A el \emph{payload} se comprime con un codificador rANS estático de
byte (renormalización de $8$ bits, total $M=2^{12}$). En el Modo~B el texto se
comprime con un esquema que combina diccionario LZ y codificación de Huffman con
modelado de contexto; sobre prosa natural corta este esquema supera a la
codificación por diccionario simple gracias a un diccionario incorporado de
términos frecuentes. En ambos casos el principio es el mismo: \emph{minimizar la
entropía representada} para reducir la carga a insertar.

\section*{4. Fundamentos de transformación y percepción}

\subsection*{4.1 Transformada Discreta del Coseno (DCT)}

Ambos modos operan en el dominio de la DCT. La DCT-II bidimensional ortonormal de
un bloque $f(x,y)$ de tamaño $N\times N$ es
\begin{equation}
F(u,v) = \alpha(u)\,\alpha(v)\sum_{x=0}^{N-1}\sum_{y=0}^{N-1} f(x,y)
\cos\!\Big[\tfrac{\pi(2x+1)u}{2N}\Big]
\cos\!\Big[\tfrac{\pi(2y+1)v}{2N}\Big],
\end{equation}
con $\alpha(0)=\sqrt{1/N}$ y $\alpha(k)=\sqrt{2/N}$ para $k>0$, lo que hace la
transformada ortonormal (energía preservada, teorema de Parseval). La DCT fue
introducida por Ahmed, Natarajan y Rao \cite{ahmed1974} y concentra la energía de
señales naturales en pocos coeficientes de baja frecuencia, propiedad que se
explota para insertar información donde es a la vez robusta e imperceptible. El
Modo~A usa bloques $8\times 8$ alineados con JPEG; el Modo~B usa una única DCT
global sobre la malla canónica.

\subsection*{4.2 Enmascaramiento perceptual del sistema visual humano}

El sistema visual humano (HVS) tolera más ruido en regiones de alta textura que en
zonas lisas o bordes definidos: es el fenómeno de \emph{enmascaramiento por
contraste}. Los modelos perceptuales sobre la DCT \cite{watson1993} formalizan
umbrales de visibilidad dependientes del contenido local. El sistema aprovecha
esto de dos maneras: (i)~en el Modo~A, la función de costo dirige las
modificaciones hacia texturas complejas; (ii)~en el Modo~B, una máscara perceptual
basada en la actividad local modula la amplitud de la marca (Sección~6.3).

\section*{5. Fundamentos de codificación de canal}

\subsection*{5.1 Códigos Reed-Solomon}

Para reparar los errores que el canal introduce pese a todo, el sistema envuelve
la carga con un código Reed-Solomon \cite{reedsolomon1960}. Un código
$\text{RS}(n,k)$ sobre el cuerpo finito $\mathrm{GF}(2^8)$ interpreta $k$ bytes de
datos como coeficientes de un polinomio y lo evalúa en $n=255$ puntos, añadiendo
$n-k$ bytes de paridad. Por la \emph{cota de distancia mínima} (distancia
$d=n-k+1$, códigos MDS \cite{macwilliams1977}) corrige hasta
\begin{equation}
t = \Big\lfloor \tfrac{n-k}{2} \Big\rfloor
\end{equation}
bytes erróneos por bloque. Esto es idóneo para \emph{ráfagas} de error: un byte
dañado cuenta como un solo símbolo con independencia de cuántos de sus bits fallen.
El Modo~A usa $n-k=120$ (corrige $60$ bytes por bloque); el Modo~B usa $n-k=32$
(corrige $16$).

\subsection*{5.2 Verificación por redundancia cíclica (CRC)}

La cabecera de cada contenedor incorpora una comprobación de redundancia cíclica
CRC-16 \cite{peterson1961}. El mensaje se interpreta como polinomio sobre
$\mathrm{GF}(2)$ y se toma el residuo módulo el polinomio generador
$g(x)=x^{16}+x^{12}+x^{5}+1$ (estándar CCITT). El CRC no corrige, pero
\emph{detecta} con altísima probabilidad cualquier corrupción de la cabecera; su
papel es evitar \emph{falsos positivos}: una imagen sin marca (o demasiado
degradada) falla el CRC y se reporta como ``sin marca'' en lugar de devolver
basura.

\section*{6. Modo A: esteganografía de alta capacidad (ANS + STC)}

Este modo persigue \emph{capacidad}: ocultar kilobytes resistiendo la
recompresión JPEG. Su tubería es
\[
\text{payload}\ \xrightarrow{\text{ANS}}\ \text{comprimido}\
\xrightarrow{\text{Reed-Solomon}}\ \text{protegido}\
\xrightarrow{\text{STC}}\ \text{coef. DCT}\ \rightarrow\ \text{imagen}.
\]

\subsection*{6.1 Inserción de paridad cuantizada en la DCT}

Sobre cada bloque $8\times 8$ se seleccionan coeficientes de \emph{media
frecuencia} (se excluyen la componente continua, muy visible, y las altas
frecuencias, que JPEG destruye primero). Un bit se incrusta en la \emph{paridad}
del coeficiente cuantizado
\begin{equation}
b = \Big\lfloor \tfrac{F(u,v)}{\Delta} \Big\rceil \bmod 2,
\end{equation}
donde $\Delta$ es un paso de cuantización elegido de modo que \emph{supere} el paso
propio de JPEG en esas frecuencias; así la paridad sobrevive a la recompresión.

\subsection*{6.2 Función de costo J-UNIWARD}

¿Qué coeficientes conviene modificar? Los de menor impacto perceptual y
estadístico. J-UNIWARD \cite{holub2014uniward} define el costo de alterar el modo
$(i,j)$ del bloque $b$ como
\begin{equation}
\rho(b,i,j) = \sum_{k=1}^{3}\sum_{u,v}
\frac{\big|\,\Delta R^{(k)}_{u,v}\,\big|}{\sigma + \big|\,R^{(k)}_{u,v}\,\big|},
\end{equation}
donde $R^{(k)}$ son los tres residuos direccionales de una descomposición
\emph{wavelet} no decimada de Daubechies-8 \cite{daubechies1992} de la imagen, y
$\Delta R^{(k)}$ es su variación ante una modificación unidad del coeficiente. La
constante $\sigma=2^{-6}$ estabiliza el cociente. En regiones lisas
$|R^{(k)}|\!\approx\!0$ y el costo se dispara (se evitan); en texturas el costo es
pequeño (se prefieren). Así el enmascaramiento perceptual se convierte en una
función objetivo cuantitativa.

\subsection*{6.3 Códigos de Síndrome-Trellis y algoritmo de Viterbi}

Dado el vector de portadores $\vect{x}$ (paridades de los coeficientes), su vector
de costos $\boldsymbol{\rho}$ y el mensaje $\vect{m}$, los Códigos de
Síndrome-Trellis (STC) \cite{filler2011} buscan el estego $\vect{y}$ que resuelve
\begin{equation}
\vect{y} = \arg\min_{\,\vect{z}\,:\;\vect H\vect{z}=\vect{m}\ (\mathrm{mod}\ 2)}\;
D(\vect{x},\vect{z}), \qquad
D(\vect{x},\vect{z}) = \sum_i \rho_i\,[\,x_i \ne z_i\,].
\end{equation}
La matriz de paridad $\vect H$ se forma repitiendo en diagonal una submatriz
compartida $\hat{\vect H}$ de altura $h$ (aquí $h=8$, con $2^h$ estados de
trellis). La minimización es un problema de \emph{camino de mínimo costo} en la
malla del síndrome, resuelto de forma óptima por el algoritmo de Viterbi
\cite{viterbi1967}. La extracción es el mapa lineal barato
$\vect{m}=\vect H\vect{y}\ (\mathrm{mod}\ 2)$.

Una consecuencia clave es la \emph{avalancha}: un solo coeficiente alterado por el
canal puede voltear hasta $h$ bits del mensaje. Precisamente por ello el envoltorio
Reed-Solomon (Sección~5.1) es indispensable: convierte esas ráfagas en errores de
byte reparables.

\section*{7. Modo B: marca de agua robusta (lo logrado)}

El Modo~A muere ante el \emph{reescalado}. El Modo~B --- el núcleo de lo logrado
--- sacrifica capacidad por robustez bruta: incrusta un texto de hasta $512$~bytes
\emph{comprimidos} y sobrevive la cadena completa de reescalado y recompresión de
las redes sociales. Sus componentes matemáticos son los siguientes.

\subsection*{7.1 Resincronización a resolución canónica}

Sea $Y$ la luminancia del portador. Tanto el insertor como el extractor la llevan
primero a una malla fija $S\times S$ (con $S=1152$) mediante remuestreo:
\[
Y_c = \mathcal{R}_{S\times S}(Y).
\]
Cualquiera que sea el tamaño arbitrario al que una plataforma reescale la imagen,
el extractor la ``devuelve'' a la misma malla canónica antes de decodificar. Con
ello la retícula de incrustación \emph{nunca se pierde}: se resuelve la
desincronización que mata a la esteganografía por bloques. La marca se calcula en
la DCT de $Y_c$ y solo la \emph{perturbación} se remuestrea de vuelta al tamaño
nativo, preservando el aspecto original.

\subsection*{7.2 Espectro ensanchado mejorado (ISS)}

La información se incrusta por Espectro Ensanchado. En el esquema aditivo clásico
\cite{cox1997} cada bit se reparte sobre muchos coeficientes con un patrón
pseudoaleatorio de signos, pero la propia imagen actúa como ruido (interferencia
del portador). El \emph{Espectro Ensanchado Mejorado} (ISS) de Malvar y Florêncio
\cite{malvar2003} cancela algebraicamente esa interferencia.

Sea $G_i$ el grupo de $L=|G_i|$ coeficientes asignados al bit $b_i$, con signos
$s_c\in\{-1,+1\}$ y amplitud $\alpha$. Definimos el objetivo
$\tau_i=\alpha(2b_i-1)$ y la proyección previa
\begin{equation}
p_i = \frac{1}{\sqrt{L}}\sum_{c\in G_i} s_c\,D(c).
\end{equation}
La regla de inserción actualiza cada coeficiente como
\begin{equation}
D'(c) = D(c) + \frac{\tau_i - p_i}{\sqrt{L}}\,s_c .
\end{equation}
Como $s_c^2=1$, la proyección \emph{posterior} resulta exactamente
\begin{equation}
p'_i = \frac{1}{\sqrt L}\sum_{c\in G_i} s_c D'(c)
     = p_i + \frac{\tau_i-p_i}{L}\sum_{c\in G_i} s_c^2 = \tau_i,
\end{equation}
es decir, la correlación se fija en $\pm\alpha$ \emph{eliminando} el término del
portador $p_i$. La detección es ciega y por \emph{signo}:
\begin{equation}
\hat{b}_i = \mathbb{1}\!\left[\ \sum_{c\in G_i} s_c\,\tilde D(c) \ > \ 0\ \right],
\end{equation}
donde $\tilde D$ es el espectro recibido tras el canal. Al no depender de la
amplitud, el extractor \emph{no necesita conocer} $\alpha$, y la ganancia de
procesamiento propia del espectro ensanchado hace la decisión muy robusta al
ruido de compresión.

\subsection*{7.3 Enmascaramiento perceptual adaptativo}

La perturbación espacial $\delta(x,y)$ (anti-transformada de la marca) se multiplica
por una máscara de actividad local para concentrarla en texturas y retirarla de
zonas lisas. Con media y varianza locales en ventana $w\times w$,
\begin{equation}
a(x,y) = \sqrt{\operatorname{Var}_w[Y_c](x,y)}, \qquad
\mathrm{mask}(x,y) = \operatorname{clip}\!\Big(\frac{a(x,y)}{\overline{a}},\ \ell,\ u\Big),
\end{equation}
con recorte a $[\ell,u]=[0{.}72,\,3{.}0]$. Este modelo de enmascaramiento por
contraste \cite{watson1993} preserva la imperceptibilidad; el suelo $\ell>0$
garantiza que las regiones lisas (cielo, piel, documentos) conserven marca
suficiente para sobrevivir descargas pequeñas.

\subsection*{7.4 Fuerza adaptativa por carga útil}

La fiabilidad de la detección por correlación crece con la \emph{ganancia de
procesamiento}: a distorsión perceptual fija, la razón señal-ruido del estadístico
de decisión aumenta aproximadamente como $\sqrt{L}$, siendo $L$ el número de
coeficientes por bit \cite{cox1997,malvar2003}. Un mensaje largo dispone de menos
coeficientes por bit ($L$ pequeño) y pierde margen. Para compensarlo, la amplitud
se escala en función de $L$:
\begin{equation}
\alpha_{\text{payload}} = \alpha \cdot \min\!\Big(\beta_{\max},\ \sqrt{\tfrac{C_{\text{ref}}}{L}}\Big),
\qquad L = \frac{N-N_{\text{cab}}}{n_{\text{bits}}},
\end{equation}
con $C_{\text{ref}}=400$ y tope $\beta_{\max}=1{.}60$. Así los mensajes cortos se
insertan con baja amplitud (alta calidad, PSNR alto) mientras que los largos suben
su amplitud $\propto 1/\sqrt{L}$ para mantener el mismo margen de canal. Como la
detección es por signo, este ajuste es \emph{transparente} para el extractor.

\subsection*{7.5 Cota de resolución de trabajo y el teorema de muestreo}

La marca ocupa frecuencias espaciales hasta $f_{\text{hi}}$ ciclos por ancho de
imagen. Por el \emph{teorema de muestreo de Nyquist-Shannon}
\cite{nyquist1928,shannon1949}, un reescalado posterior a $W$ píxeles solo preserva
frecuencias por debajo de la frecuencia de Nyquist $W/2$: todo componente por
encima es aniquilado por el filtro anti-aliasing. Insertar sobre un original muy
grande obliga a la marca a atravesar un remuestreo de gran factor
(canónico~$\to$~nativo al insertar, y nativo~$\to$~plataforma al descargar); la
composición de dos núcleos de interpolación imperfectos atenúa la banda media y
hace fallar la recuperación de forma \emph{dependiente de la razón de remuestreo}.

La solución adoptada acota el tamaño de trabajo: todo original cuyo lado mayor
exceda $S_{\max}=2048$ se reduce a $S_{\max}$ antes de insertar. Esto (i)~acota la
razón de remuestreo canónico-nativo-descarga, y (ii)~coincide con el propio límite
que las plataformas aplican a las imágenes. Es una aplicación directa del teorema
de muestreo y constituyó la corrección clave para el fallo observado en imágenes
de alta resolución.

\subsection*{7.6 Contenedor: compresión, Reed-Solomon y cabecera}

La carga final es
\[
\text{texto}\ \xrightarrow{\text{cod. entrópica}}\ \text{comprimido}\
\xrightarrow{\text{Reed-Solomon}}\ \text{región de payload},
\]
precedida de una cabecera ultra-redundante de $64$~bits
$\;\texttt{MAGIC}\,|\,\texttt{LEN}\,|\,\texttt{NSYM}\,|\,\texttt{CRC-16}\,|\,\texttt{FLAGS}$.
La cabecera se decodifica primero y anuncia cuántos bits de payload leer; los
mensajes cortos usan menos bits (mayor calidad y margen). El bit de compresión
viaja en \texttt{FLAGS}; como la comprobación CRC cubre solo
$\texttt{MAGIC}|\texttt{LEN}|\texttt{NSYM}$, las marcas de versiones previas
(cuyo octavo byte era relleno nulo) se interpretan como no comprimidas y siguen
verificándose (compatibilidad hacia atrás). La compresión se aplica solo si
reduce el tamaño; cadenas cortas o incompresibles se almacenan en crudo. En la
decodificación se acota el tamaño expandido para prevenir un \emph{ataque de
descompresión} (bomba).

\section*{8. Métricas de validación}

\begin{table}[H]
\centering
\renewcommand{\arraystretch}{1.4}
\begin{tabular}{@{}p{3.2cm}p{2.4cm}p{7.6cm}@{}}
\toprule
\textbf{Métrica} & \textbf{Dominio} & \textbf{Objetivo} \\
\midrule
PSNR & Visual & Garantizar imperceptibilidad. \\
SSIM & Estructural & Preservar la composición de la imagen. \\
BER & Integridad & Minimizar errores tras el canal. \\
Divergencia KL & Estadístico & Resistir el esteganálisis moderno. \\
\bottomrule
\end{tabular}
\caption{Métricas y su propósito en el proyecto.}
\end{table}

\noindent\textbf{PSNR.} La relación señal-ruido de pico, sobre valores en
$[0,255]$, es
\begin{equation}
\mathrm{PSNR} = 10\,\log_{10}\!\frac{255^2}{\mathrm{MSE}}, \qquad
\mathrm{MSE}=\frac{1}{MN}\sum_{x,y}\big(I(x,y)-\hat I(x,y)\big)^2 .
\end{equation}

\noindent\textbf{SSIM.} El índice de similitud estructural \cite{wang2004ssim}
compara luminancia, contraste y estructura en ventanas locales,
\begin{equation}
\mathrm{SSIM}(a,b)=\frac{(2\mu_a\mu_b+c_1)(2\sigma_{ab}+c_2)}
{(\mu_a^2+\mu_b^2+c_1)(\sigma_a^2+\sigma_b^2+c_2)},
\end{equation}
más acorde con la percepción humana que el error cuadrático.

\noindent\textbf{BER.} La tasa de error de bit es la fracción de bits del mensaje
mal recuperados; en el sistema debe caer bajo la capacidad correctora del código
Reed-Solomon para que la reconstrucción sea exacta.

\noindent\textbf{Divergencia de Kullback-Leibler.} La seguridad esteganográfica se
mide por la divergencia entre la distribución del portador $P$ y la del estego $Q$
\cite{kullback1951}:
\begin{equation}
D_{\mathrm{KL}}(P\Vert Q) = \sum_{x} P(x)\,\log\frac{P(x)}{Q(x)} .
\end{equation}
Cachin \cite{cachin1998} formaliza la $\varepsilon$-seguridad de un esquema como
$D_{\mathrm{KL}}(P\Vert Q)\le\varepsilon$; minimizar el número de cambios (el
objetivo explícito de STC) reduce esta divergencia y, con ella, la detectabilidad
frente a esteganálisis por modelos ricos y clasificadores de conjunto
\cite{fridrich2012srm,kodovsky2012ensemble}.

\section*{9. Resultados obtenidos}

La marca de agua robusta implementada:
\begin{itemize}
\item \textbf{Sobrevive la cadena social completa} (reescalado por lado mayor +
recompresión JPEG/WebP con submuestreo de croma) hasta descargas de $480$~px, para
textos cortos y para párrafos, sobre portadores texturados y suaves.
\item \textbf{Corrige el fallo de alta resolución} mediante la cota de resolución
de trabajo (Sección~7.5): originales de $3000$--$6000$~px, acotados a $2048$~px,
recuperan la marca en todo el rango de descargas.
\item \textbf{Reduce la entropía del mensaje} en torno al $45$--$48\%$ sobre prosa
natural (p.\,ej.\ un párrafo de $\sim$250 caracteres baja a $\sim$130~bytes),
elevando el margen de canal y la calidad de imagen.
\item \textbf{Preserva la imperceptibilidad}, con PSNR de $\sim$38~dB para mensajes
cortos y $\sim$31~dB para cargas cercanas al máximo, y decae con gracia (reporta
``sin marca'' en lugar de basura) cuando no hay marca válida.
\end{itemize}
Estos resultados confirman la premisa central: \emph{minimizar el número de
cambios en el portador} --- vía compresión entrópica, costo adaptativo y
codificación de mínima distorsión --- maximiza a la vez la robustez y la
imperceptibilidad.

\section*{10. Conclusiones}

Se construyó un sistema que resuelve el problema de \emph{persistencia de marcas de
agua invisibles} bajo canales con pérdida combinando, de forma coherente,
resultados clásicos de teoría de la información (entropía de Shannon, codificación
de Huffman/LZ/ANS), procesamiento de señales (DCT, espectro ensanchado mejorado,
teorema de muestreo), codificación de canal (Reed-Solomon, CRC) y optimización de
mínima distorsión (funciones de costo UNIWARD, códigos Síndrome-Trellis con
Viterbi). El Modo~B, basado en ISS sobre una malla canónica con amplitud adaptativa
y cota de resolución, logra el objetivo práctico de sobrevivir a Facebook, WhatsApp
y Pinterest sin sacrificar la invisibilidad, aportando una base matemática sólida
para una marca de procedencia anti-\emph{deepfake}.

\begin{thebibliography}{99}

\bibitem{shannon1948}
C.~E. Shannon, ``A Mathematical Theory of Communication,''
\emph{Bell System Technical Journal}, vol.~27, pp.~379--423 y 623--656, 1948.

\bibitem{huffman1952}
D.~A. Huffman, ``A Method for the Construction of Minimum-Redundancy Codes,''
\emph{Proceedings of the IRE}, vol.~40, no.~9, pp.~1098--1101, 1952.

\bibitem{zivlempel1977}
J.~Ziv y A.~Lempel, ``A Universal Algorithm for Sequential Data Compression,''
\emph{IEEE Transactions on Information Theory}, vol.~23, no.~3, pp.~337--343, 1977.

\bibitem{rissanen1979}
J.~Rissanen y G.~G. Langdon, ``Arithmetic Coding,''
\emph{IBM Journal of Research and Development}, vol.~23, no.~2, pp.~149--162, 1979.

\bibitem{duda2014}
J.~Duda, ``Asymmetric numeral systems: entropy coding combining speed of Huffman
with compression rate of arithmetic coding,''
\emph{arXiv:1311.2540}, 2014.

\bibitem{ahmed1974}
N.~Ahmed, T.~Natarajan y K.~R. Rao, ``Discrete Cosine Transform,''
\emph{IEEE Transactions on Computers}, vol.~C-23, no.~1, pp.~90--93, 1974.

\bibitem{daubechies1992}
I.~Daubechies, \emph{Ten Lectures on Wavelets}. Philadelphia: SIAM, 1992.

\bibitem{watson1993}
A.~B. Watson, ``DCT quantization matrices visually optimized for individual
images,'' en \emph{Proc.\ SPIE Human Vision, Visual Processing, and Digital
Display IV}, vol.~1913, pp.~202--216, 1993.

\bibitem{reedsolomon1960}
I.~S. Reed y G.~Solomon, ``Polynomial Codes over Certain Finite Fields,''
\emph{Journal of the Society for Industrial and Applied Mathematics}, vol.~8,
no.~2, pp.~300--304, 1960.

\bibitem{macwilliams1977}
F.~J. MacWilliams y N.~J.~A. Sloane, \emph{The Theory of Error-Correcting Codes}.
Amsterdam: North-Holland, 1977.

\bibitem{peterson1961}
W.~W. Peterson y D.~T. Brown, ``Cyclic Codes for Error Detection,''
\emph{Proceedings of the IRE}, vol.~49, no.~1, pp.~228--235, 1961.

\bibitem{cox1997}
I.~J. Cox, J.~Kilian, F.~T. Leighton y T.~Shamoon, ``Secure Spread Spectrum
Watermarking for Multimedia,'' \emph{IEEE Transactions on Image Processing},
vol.~6, no.~12, pp.~1673--1687, 1997.

\bibitem{malvar2003}
H.~S. Malvar y D.~A.~F. Florêncio, ``Improved Spread Spectrum: A New Modulation
Technique for Robust Watermarking,'' \emph{IEEE Transactions on Signal
Processing}, vol.~51, no.~4, pp.~898--905, 2003.

\bibitem{filler2011}
T.~Filler, J.~Judas y J.~Fridrich, ``Minimizing Additive Distortion in
Steganography Using Syndrome-Trellis Codes,'' \emph{IEEE Transactions on
Information Forensics and Security}, vol.~6, no.~3, pp.~920--935, 2011.

\bibitem{viterbi1967}
A.~J. Viterbi, ``Error Bounds for Convolutional Codes and an Asymptotically
Optimum Decoding Algorithm,'' \emph{IEEE Transactions on Information Theory},
vol.~13, no.~2, pp.~260--269, 1967.

\bibitem{li2014hill}
B.~Li, M.~Wang, J.~Huang y X.~Li, ``A New Cost Function for Spatial Image
Steganography,'' en \emph{Proc.\ IEEE International Conference on Image Processing
(ICIP)}, pp.~4206--4210, 2014.

\bibitem{holub2014uniward}
V.~Holub, J.~Fridrich y T.~Denemark, ``Universal Distortion Function for
Steganography in an Arbitrary Domain,'' \emph{EURASIP Journal on Information
Security}, vol.~2014, no.~1, art.~1, 2014.

\bibitem{nyquist1928}
H.~Nyquist, ``Certain Topics in Telegraph Transmission Theory,''
\emph{Transactions of the AIEE}, vol.~47, no.~2, pp.~617--644, 1928.

\bibitem{shannon1949}
C.~E. Shannon, ``Communication in the Presence of Noise,''
\emph{Proceedings of the IRE}, vol.~37, no.~1, pp.~10--21, 1949.

\bibitem{wang2004ssim}
Z.~Wang, A.~C. Bovik, H.~R. Sheikh y E.~P. Simoncelli, ``Image Quality Assessment:
From Error Visibility to Structural Similarity,'' \emph{IEEE Transactions on Image
Processing}, vol.~13, no.~4, pp.~600--612, 2004.

\bibitem{kullback1951}
S.~Kullback y R.~A. Leibler, ``On Information and Sufficiency,''
\emph{The Annals of Mathematical Statistics}, vol.~22, no.~1, pp.~79--86, 1951.

\bibitem{cachin1998}
C.~Cachin, ``An Information-Theoretic Model for Steganography,'' en
\emph{Information Hiding}, LNCS 1525, pp.~306--318, 1998.

\bibitem{fridrich2012srm}
J.~Fridrich y J.~Kodovský, ``Rich Models for Steganalysis of Digital Images,''
\emph{IEEE Transactions on Information Forensics and Security}, vol.~7, no.~3,
pp.~868--882, 2012.

\bibitem{kodovsky2012ensemble}
J.~Kodovský, J.~Fridrich y V.~Holub, ``Ensemble Classifiers for Steganalysis of
Digital Media,'' \emph{IEEE Transactions on Information Forensics and Security},
vol.~7, no.~2, pp.~432--444, 2012.

\end{thebibliography}

\end{document}
